\documentclass{article}
\usepackage{amsmath,amsfonts,amssymb}
\usepackage{graphicx}
\usepackage{hyperref}

\title{Hierarchical Skill-Based Navigation with Reinforcement Learning for Temporal Knowledge Graph Question Answering}
\author{TKGQA Navigation Agent System}
\date{}

\begin{document}

\maketitle

\begin{abstract}
We propose a hierarchical navigation framework for Temporal Knowledge Graph Question Answering (TKGQA), where reasoning is formulated as a learned navigation problem over a structured skill space. Our system integrates skill abstraction, neural routing, and reinforcement learning to enable efficient multi-hop reasoning over temporal knowledge graphs. Unlike prior work relying on flat retrieval or pure LLM reasoning, our method constructs a navigable skill tree and optimizes navigation policies via trajectory-level rewards.
\end{abstract}

\section{Introduction}
Temporal Knowledge Graph Question Answering (TKGQA) requires reasoning over entities, relations, and temporal constraints. Existing approaches suffer from flat retrieval inefficiency and lack structured navigation. We address this by introducing a hierarchical skill-based navigation framework.

\section{Method}

\subsection{Skill Abstraction}
We define a skill as a structured unit consisting of entity sets, relation types, temporal scope, and associated document identifiers.

\subsection{Hierarchical Navigation}
The system performs multi-step navigation:
Query $\rightarrow$ Skill Router $\rightarrow$ Neural Policy $\rightarrow$ Navigator $\rightarrow$ Document Retrieval.

\subsection{Neural Policy}
We model skill selection as:
\begin{equation}
\pi(s|q) = \frac{\exp(\cos(e_q, e_s))}{\sum_{s'} \exp(\cos(e_q, e_{s'}))}
\end{equation}

\subsection{Reinforcement Learning Optimization}
We optimize the policy using trajectory-level rewards:
\begin{equation}
\nabla J(\theta) = \mathbb{E}[\log \pi(a|s)(R - b)]
\end{equation}
where $b$ is a baseline.

\section{Unified Agent}
We unify routing, navigation, neural policy, and optimization into a single agent loop that performs both inference and learning.

\section{Experiments}

\subsection{Baselines}
Flat retrieval, heuristic routing, and LLM-only reasoning.

\subsection{Metrics}
Accuracy, Hit@K, navigation depth, and retrieval efficiency.

\subsection{Ablation}
We remove neural policy, RL optimization, and hierarchical structure to evaluate contributions.

\section{Conclusion}
We present a unified hierarchical navigation framework for TKGQA that integrates structured decomposition, neural ranking, and reinforcement learning into a single agent system.

\end{document}
